\documentclass[11pt,a4paper]{article}
\usepackage[T1]{fontenc}
\usepackage{lmodern}
\usepackage[margin=27mm,headheight=14pt]{geometry}
\usepackage{amsmath,amssymb,amsthm,mathtools,bm}
\usepackage{microtype,booktabs,longtable,array,enumitem}
\usepackage{xcolor,hyperref,fancyhdr,needspace}
\hypersetup{colorlinks=true,linkcolor=blue!45!black,citecolor=blue!45!black,urlcolor=blue!45!black,pdftitle={Ordered Hurwitz Germs at Every Depth},pdfauthor={ProveIt research continuation}}
\pagestyle{fancy}\fancyhf{}
\fancyhead[L]{\small Ordered Hurwitz germs}\fancyhead[R]{\small ProveIt research continuation}
\fancyfoot[C]{\thepage}
\setlength{\parskip}{4pt}\setlength{\parindent}{1.3em}
\setlist{nosep,leftmargin=*}
\numberwithin{equation}{section}
\newtheorem{theorem}{Theorem}[section]
\newtheorem{lemma}[theorem]{Lemma}
\newtheorem{corollary}[theorem]{Corollary}
\newtheorem{proposition}[theorem]{Proposition}
\theoremstyle{definition}\newtheorem{definition}[theorem]{Definition}
\theoremstyle{remark}\newtheorem{remark}[theorem]{Remark}
\DeclareMathOperator{\Li}{Li}
\DeclareMathOperator{\FP}{FP}
\DeclareMathOperator{\CT}{CT}
\DeclareMathOperator{\Res}{Res}
\DeclareMathOperator{\ord}{ord}
\newcommand{\eps}{\varepsilon}
\newcommand{\HH}{\mathcal H}
\newcommand{\Atail}{\mathcal A}
\newcommand{\LL}{\mathcal L}
\newcommand{\PP}{\mathcal P}
\newcommand{\RR}{\mathcal R}
\newcommand{\dd}{\,\mathrm d}
\newcommand{\src}[1]{{\small\nolinkurl{#1}}}
\newcommand{\ones}{\mathbf 1}
\newcommand{\gammak}[1]{\gamma_{#1}(a)}

\title{\textbf{Ordered Hurwitz Germs at Every Depth}\\[5pt]
\large Harmonic renormalization, directional Stieltjes constants,\\
and polylogarithmic Mellin identities}
\author{A research continuation prepared for the ProveIt project}
\date{10 October 2026}

\begin{document}
\maketitle
\begin{abstract}
We construct a compatible hierarchy of holomorphic remainders for ordered
multiple Hurwitz zeta functions near the all-one point. For prefix sums
$U_k=u_1+\cdots+u_k$, the exact decomposition is
\[
 Z_d(1+u_1,\ldots,1+u_d;a)
 =\sum_{j=0}^{d}
 \frac{\HH_j(u_{d-j+1},\ldots,u_d;a)}{U_1\cdots U_{d-j}},
 \qquad \HH_0=1.
\]
A normally convergent outer-tail subtraction and a divided-difference
recursion construct every remainder and every mixed Taylor coefficient.
We identify these coefficients with constant terms of logarithmic harmonic
sums. This gives a specified stuffle-preserving regularization and explicit
counterterms converting arbitrary transverse directional constants to it.
At depth three we obtain a complete formula for every ordered ray; its
nonsymmetric dependence is carried by one explicitly normalized convergent
Stieltjes sum. Nonlinear and tangential curves are also treated, with the
precise finite jet of the curve needed for a constant term.

The regular all-one values are generated by $\Gamma(a)/\Gamma(a+z)$.
The coefficient recursion yields ordinary convergent integrals involving
multiple polylogarithms and their order derivatives. At $a=1$, an all-depth
family reduces to elementary logarithmic Mellin integrals and has the
closed generating identity
\[
 \int_0^\infty\left(\frac1{e^t-1}-\frac1t\right)t^z
       \bigl((1-e^{-t})^{-z}-1\bigr)\dd t
 =\frac1z-\Gamma(1+z)\zeta(1+z),\qquad \Re z>-1,
\]
with a removable value at zero. We give a Hurwitz-shifted hypergeometric
counterpart, rational-shift cyclotomic coordinates, and exact polygamma
primitives. The article resolves the ordered-germ and specified-subtraction
targets recorded in the inspected repository continuation; it does not
claim priority for classical multiple-zeta continuation or regularization,
and it does not prove the Gaussian $S_6$ or current $S_8$ conjectures.
\end{abstract}

\clearpage
\tableofcontents
\newpage

\section{Source target, contribution, and mathematical status}
\label{sec:scope}

The immediate source is the ProveIt continuation
\emph{resonant-jets-and-shifts}, especially
\src{sections/03-directional-hurwitz.tex} and
\src{sections/08-audit-research.tex}. The former proves the exact ordered
double-Hurwitz germ and its directional constants. The latter explicitly
asks for an ordered depth-three analogue, compatible local subtraction,
and tangential approaches to nested poles. Its question 6 asks for a useful
formula for the nonsymmetric Stieltjes sums. These are concrete continuation
targets, not assertions that the corresponding general subjects are absent
from the literature \cite{repo-target,repo-double}.

The results below settle the following specified tasks. Theorem
\ref{thm:germ} constructs the ordered germ at every finite depth, not merely
a fully symmetrized value. Theorem \ref{thm:cutoff} identifies its regular
coefficients with harmonic cutoff constants, and Theorem
\ref{thm:stuffle} proves the associated stuffle property. Corollary
\ref{cor:counterterm} supplies the exact directional counterterms. Theorem
\ref{thm:tripleFP} evaluates every transverse ordered depth-three ray.
Theorem \ref{thm:curve} treats every analytic curve through the all-one
point that is not contained in a polar divisor. Proposition
\ref{prop:omega} supplies a convergent formula for the first nonsymmetric
coordinate, without claiming that it reduces to ordinary one-variable
constants.

The external comparison is important. Matsumoto, Onozuka and Wakabayashi
already give an all-depth analytic polar decomposition and Laurent
coefficient algorithms for Euler--Zagier zeta functions \cite{mow}.
Their ordering is opposite to ours, and their single-variable constants
absorb the sign and factorial in the usual Stieltjes convention. The
present construction is a common-Hurwitz-shift refinement with a particular
\emph{harmonic-compatible} holomorphic remainder, an explicit ordinary
Mellin realization, and the stated directional and curve formulas.
General meromorphic continuation, harmonic products, and the existence of
stuffle-preserving regularizations are not claimed as new principles.
Guo--Zhang and the related renormalization literature already establish
such principles in much broader settings \cite{gz,gpxz}.

\paragraph{Audit boundary.}
The manuscript README, the incoming directory listing, and the named
textual continuation were inspected. The repository commit recorded in the
provenance manifest is
\src{0c740647c48baec7f0b4192723e977d69917dd52}; the two decisive source
files have independently recorded Git blob identifiers. The six ZIP
archives then visible in \src{docs/incoming} were inventoried by name and
metadata, not exhaustively unpacked or re-proved. Accordingly, claims of
progress are relative to the explicit textual targets just named, not a
claim of exhaustive novelty across every incoming archive. The current
manuscript distinguishes proved statements, preserved reports, and
conjectural arithmetic evaluations; that distinction is retained
\cite{repo-readme,repo-incoming}.

\paragraph{Proof status.}
All numbered mathematical results in this article are proved here under
their stated hypotheses. Finite symbolic and floating-point checks are
separate diagnostics. No proof-assistant formalization, interval enclosure
of the numerical tests, arithmetic independence theorem, or exhaustive
literature-priority claim is made.

\section{Conventions and the ordered multiple Hurwitz function}
\label{sec:notation}

Fix $a>0$. The larger summation index always carries the earlier exponent:
\begin{equation}\label{eq:Zd}
 Z_d(s_1,\ldots,s_d;a)
 =\sum_{n_1>\cdots>n_d\ge0}\prod_{j=1}^{d}(n_j+a)^{-s_j},
 \qquad Z_0=1.
\end{equation}
At $s_j=1+u_j$, the absolute-convergence domain is
$\Re U_k>0$ for $1\le k\le d$, where
\begin{equation}\label{eq:prefix}
 U_k=u_1+\cdots+u_k,\qquad U_0=0.
\end{equation}
We may first prove identities on the smaller domain $\Re u_j>0$ and then
continue them as meromorphic germs.

Our Stieltjes normalization is
\begin{equation}\label{eq:g}
 \zeta(1+u,a)=\frac1u+g_a(u),\qquad
 g_a(u)=\sum_{m\ge0}g_m(a)u^m,
 \qquad g_m(a)=\frac{(-1)^m\gamma_m(a)}{m!}.
\end{equation}
In particular $\gamma_0(a)=-\psi(a)$. A prime on $\zeta(s,a)$ means
an order derivative in $s$; a prime on a function of $a$ means an ordinary
$a$-derivative. Standard Hurwitz recurrences and the relation to
polygamma functions are as in DLMF \cite{dlmf-zeta}.

For $\bm m=(m_1,\ldots,m_d)\in\mathbb Z_{\ge0}^d$, write
$|\bm m|=\sum m_i$ and $\bm m!=\prod m_i!$. Coefficients of a holomorphic
function are denoted by $[\bm u^{\bm m}]$. A directional finite part is
\emph{always} the coefficient $[\eps^0]$ after the displayed substitution.
It is not a path-independent multivariate value.

We use the elementary-subtracted Bose kernel
\begin{equation}\label{eq:kernel}
 K(t)=\frac1{e^t-1}-\frac1t,\qquad t>0.
\end{equation}
It is bounded at zero, with $K(0)=-1/2$, and is $O(1/t)$ at infinity.
These estimates are used only to justify identities and differentiations,
not as the principal research conclusions.

\section{A normally convergent outer-tail subtraction}
\label{sec:tail}

\begin{lemma}[One-variable tail]\label{lem:tail}
For $x>0$, define
\begin{equation}\label{eq:B}
 B(u,x)=\zeta(1+u,x+1)-\frac{x^{-u}}u,
\end{equation}
where the singularity at $u=0$ is removed. Then $B$ is entire in $u$, and
uniformly on compact subsets of $\Re u>-1$,
\begin{equation}\label{eq:Bint}
 B(u,x)=\frac1{\Gamma(1+u)}
       \int_0^\infty t^u e^{-xt}K(t)\dd t.
\end{equation}
As $x\to\infty$, locally uniformly in $u$ near zero,
\begin{equation}\label{eq:Basym}
 B(u,x)=-\frac12x^{-1-u}+O(x^{-2-\Re u}).
\end{equation}
Fixed $u$-derivatives obey the corresponding estimate with a fixed power
of $1+\log x$.
\end{lemma}
\begin{proof}
Both terms in \eqref{eq:B} have residue one at $u=0$, so their difference
is entire. For $\Re u>0$, subtract the standard Mellin integral for
$x^{-u}/u$ from that for $\zeta(1+u,x+1)$; this gives
\eqref{eq:Bint}. Since $K(t)=-1/2+O(t)$ at zero and the exponential handles
infinity, the integral is holomorphic for $\Re u>-1$ and provides the
continuation there. Substituting $t=y/x$ proves \eqref{eq:Basym} by the
same expansion. Differentiation under the integral, or Cauchy estimates
on slightly larger compact $u$-sets, gives the derivative assertion.
\end{proof}

\begin{definition}[Holomorphic tail terms]\label{def:A}
For $d\ge2$, set
\begin{equation}\label{eq:Ad}
 \Atail_d(\bm u;a)=
 \sum_{n_2>\cdots>n_d\ge0}
 B(u_1,n_2+a)\prod_{j=2}^{d}(n_j+a)^{-1-u_j}.
\end{equation}
\end{definition}

\begin{lemma}[Normal convergence and outer-tail recursion]\label{lem:normal}
The series \eqref{eq:Ad}, together with every fixed mixed derivative,
converges normally near $\bm u=0$. In particular it defines a holomorphic
function on the neighborhood
$\sum_j|u_j|<1/4$. As meromorphic germs,
\begin{equation}\label{eq:outer}
 Z_d(1+u_1,\ldots,1+u_d;a)
 =\frac1{u_1}Z_{d-1}(1+u_1+u_2,1+u_3,\ldots,1+u_d;a)
    +\Atail_d(\bm u;a).
\end{equation}
\end{lemma}
\begin{proof}
For the normal-convergence claim, put $x=n_2+a$. On a compact set with
$\sum|u_j|\le\sigma<1/4$, the inner $(d-2)$-fold finite sum is bounded by
$C x^{\sum_{j\ge3}|u_j|}(1+\log x)^{d-2}$ for large $x$. Lemma
\ref{lem:tail} therefore bounds the grouped summand by
$C x^{-2+\sigma}(1+\log x)^{d-2}$. This is summable. Derivatives add
only fixed logarithmic powers, or follow by Cauchy estimates after a
small enlargement of the compact set. Small $x$ contributes only finitely
many terms.

In the initial convergence domain, sum the outer index $n_1>n_2$ first.
Its sum is $\zeta(1+u_1,n_2+a+1)$. Substitution of \eqref{eq:B} gives
\eqref{eq:outer}. The normally convergent remainder then proves the
identity of meromorphic germs.
\end{proof}

\section{The compatible polar germ at every depth}
\label{sec:germ}

Define $\HH_0=1$, $\HH_1(u;a)=g_a(u)$, and recursively, for $d\ge2$,
\begin{equation}\label{eq:Hrec}
 \boxed{\displaystyle
 \HH_d(\bm u;a)=\Atail_d(\bm u;a)+
 \frac{\HH_{d-1}(u_1+u_2,u_3,\ldots,u_d;a)
       -\HH_{d-1}(u_2,u_3,\ldots,u_d;a)}{u_1}.}
\end{equation}
The quotient is interpreted by its removable value at $u_1=0$.
Equivalently, it is the ordinary integral
\begin{equation}\label{eq:DDint}
 \int_0^1\partial_1\HH_{d-1}(u_2+t u_1,u_3,\ldots,u_d;a)\dd t.
\end{equation}
The $\ell^1$ neighborhood in Lemma \ref{lem:normal} is convenient because
merging $u_1$ and $u_2$ does not increase $\sum|u_j|$.

\begin{theorem}[All-depth ordered Hurwitz germ]\label{thm:germ}
The functions \eqref{eq:Hrec} are holomorphic near the origin and satisfy
\begin{equation}\label{eq:germ}
 \boxed{\displaystyle
 Z_d(1+u_1,\ldots,1+u_d;a)
 =\sum_{j=0}^{d}
   \frac{\HH_j(u_{d-j+1},\ldots,u_d;a)}{U_1\cdots U_{d-j}}.}
\end{equation}
Empty products and $\HH_0$ are one. This is the unique compatible
hierarchy with $\HH_0=1$ and $\HH_1=g_a$ satisfying
\eqref{eq:germ} at every depth. The local polar divisors are precisely
$U_1=0,\ldots,U_d=0$. They are normal crossings in the independent linear
coordinates $U_1,\ldots,U_d$.
\end{theorem}
\begin{proof}
Holomorphy follows inductively from Lemma \ref{lem:normal} and the
removable divided difference \eqref{eq:DDint}. The claim at depth one is
\eqref{eq:g}. Suppose the claim holds at depth $d-1$. Insert that claim,
with first variable $u_1+u_2$, into \eqref{eq:outer}. Every term except
its final holomorphic term already has the required prefix denominator
and unchanged suffix numerator. The remaining expression is
\[
 \frac{\HH_{d-1}(u_1+u_2,u_3,\ldots,u_d;a)}{u_1}
   +\Atail_d(\bm u;a).
\]
Add and subtract $\HH_{d-1}(u_2,u_3,\ldots,u_d;a)/u_1$ and use
\eqref{eq:Hrec}. This proves \eqref{eq:germ}. Once all lower depths are
fixed, its depth-$d$ holomorphic term is uniquely determined by subtraction.

The displayed denominators show that there are no other local polar
divisors. Each listed divisor actually occurs: its residue, computed
below, is a nonzero meromorphic function. The invertible change of
coordinates $u_j=U_j-U_{j-1}$ gives the normal-crossing assertion.
\end{proof}

\begin{corollary}[Residues along every prefix divisor]\label{cor:residue}
For $1\le k\le d$, with the remaining variables restricted to $U_k=0$,
\begin{equation}\label{eq:residue}
 \Res_{U_k=0}Z_d
 =\frac{Z_{d-k}(1+u_{k+1},\ldots,1+u_d;a)}{U_1\cdots U_{k-1}}.
\end{equation}
For $k=d$, the numerator is $Z_0=1$.
\end{corollary}
\begin{proof}
Multiply \eqref{eq:germ} by $U_k$ and restrict to the divisor. Terms
whose denominator does not contain $U_k$ vanish; the other terms are
exactly the lower-depth germ in the numerator of \eqref{eq:residue}.
\end{proof}

\subsection{A finite coefficient recursion}
Write
\begin{equation}\label{eq:hcoeffdef}
 \HH_d(\bm u;a)=\sum_{\bm m\ge0}h_{d;\bm m}(a)\bm u^{\bm m},
 \qquad A_{d;\bm m}(a)=[\bm u^{\bm m}]\Atail_d(\bm u;a).
\end{equation}

\begin{theorem}[All ordered regular jets]\label{thm:coefficients}
For $d\ge2$ and every nonnegative multi-index,
\begin{equation}\label{eq:hrec}
 \boxed{\displaystyle
 h_{d;(m_1,\ldots,m_d)}
 =A_{d;(m_1,\ldots,m_d)}+
 \binom{m_1+m_2+1}{m_1+1}
 h_{d-1;(m_1+m_2+1,m_3,\ldots,m_d)}.}
\end{equation}
The recursion ends at $h_{1;(m)}=g_m(a)$. It uses at most $d-1$
depth reductions, and its tail coefficients are normally convergent
series or ordinary Mellin integrals.
\end{theorem}
\begin{proof}
In the divided difference, a monomial of degree
$m_1+m_2+1$ in the first argument contributes
\[
 [u_1^{m_1}u_2^{m_2}]
 \frac{(u_1+u_2)^{m_1+m_2+1}-u_2^{m_1+m_2+1}}{u_1}
 =\binom{m_1+m_2+1}{m_1+1}.
\]
The other variables are unchanged. Lemma \ref{lem:normal} justifies
coefficient extraction in $\Atail_d$.
\end{proof}

This is an algorithm for individual ordered coefficients. It does not
replace them by their full permutation average, and it does not claim a
finite basis consisting only of ordinary single-variable zeta constants.

\section{Gamma generation of the regular diagonal}
\label{sec:gamma}

Put
\begin{equation}\label{eq:Cdef}
 C_d(a)=\HH_d(0,\ldots,0;a),\qquad C_0(a)=1.
\end{equation}

\begin{theorem}[Exact diagonal generating series]\label{thm:gamma}
As an identity of formal series in $z$, with coefficients holomorphic in
$t$ near zero,
\begin{align}\label{eq:diaggen}
 \sum_{d\ge0}\HH_d(t,\ldots,t;a)z^d
 =\exp\left\{zg_a(t)+
 \sum_{k\ge2}\frac{(-1)^{k+1}}k\zeta(k+kt,a)z^k\right\}.
\end{align}
Consequently,
\begin{equation}\label{eq:gamma}
 \boxed{\displaystyle \sum_{d\ge0}C_d(a)z^d
                 =\frac{\Gamma(a)}{\Gamma(a+z)}.}
\end{equation}
\end{theorem}
\begin{proof}
The finite elementary-symmetric identity, followed by an absolutely
convergent limit for $t>0$, gives
\[
 \sum_{d\ge0}Z_d(1+t,\ldots,1+t;a)z^d
 =\exp\left\{\sum_{k\ge1}\frac{(-1)^{k+1}}k
                         \zeta(k+kt,a)z^k\right\}.
\]
On the other hand, the prefix denominators in \eqref{eq:germ} are
$(d-j)!t^{d-j}$ on the diagonal. Thus the left side is
$e^{z/t}\sum_d\HH_d(t,\ldots,t;a)z^d$. Cancel $e^{z/t}$ coefficientwise
to obtain \eqref{eq:diaggen}, and continue its coefficients to zero.
At $t=0$, the logarithm of the right side is
\[
 -\psi(a)z+\sum_{k\ge2}\frac{(-1)^{k+1}}k\zeta(k,a)z^k
 =\log\Gamma(a)-\log\Gamma(a+z),
\]
by the usual Gamma and polygamma Taylor identities \cite{dlmf-gamma,dlmf-zeta}.
This proves \eqref{eq:gamma}.
\end{proof}

With $G=\gamma_0(a)$ and $Z_k=\zeta(k,a)$, the first values are
\begin{align}
 C_1&=G,\qquad C_2=\frac{G^2-Z_2}{2},\label{eq:C2}\\
 C_3&=\frac{G^3-3GZ_2+2Z_3}{6},\label{eq:C3}\\
 C_4&=\frac{G^4-6G^2Z_2+3Z_2^2+8GZ_3-6Z_4}{24}.
 \label{eq:C4}
\end{align}
Every $C_d$ is a finite polygamma polynomial. Every \emph{diagonal}
regular derivative in \eqref{eq:diaggen} also reduces to ordinary
Stieltjes constants and order derivatives of Hurwitz zeta at positive
integers. The individual mixed coefficients in Theorem
\ref{thm:coefficients} need not collapse in this way.

\section{Harmonic cutoff constants and a compatible subtraction algebra}
\label{sec:cutoff}

Define the finite ordered sums
\begin{equation}\label{eq:FN}
 F_{d,N}(\bm u;a)=
 \sum_{N>n_1>\cdots>n_d\ge0}\prod_{j=1}^{d}(n_j+a)^{-1-u_j},
 \qquad L_N=\log(N+a).
\end{equation}
Let $\PP_0=1$ and define, with $L$ real initially,
\begin{equation}\label{eq:Pcut}
 \PP_d(\bm u;L,a)=\HH_d(\bm u;a)
   +\int_0^L e^{-u_1t}\PP_{d-1}(u_2,\ldots,u_d;t,a)\dd t.
\end{equation}
These are entire in $L$ and locally holomorphic in $\bm u$.

\begin{theorem}[The regular germ is the harmonic constant term]
\label{thm:cutoff}
Locally uniformly in $\bm u$ near zero,
\begin{equation}\label{eq:cutasym}
 F_{d,N}(\bm u;a)=\PP_d(\bm u;L_N,a)+o(1)
 \quad(N\to\infty),
\end{equation}
with an error tending to zero after every fixed mixed derivative in
$\bm u$. For each multi-index $\bm m$, the coefficient
$[\bm u^{\bm m}]\PP_d$ is a polynomial in $L$ of degree at most
$d+|\bm m|$, and its constant term is $h_{d;\bm m}(a)$. Thus
\begin{equation}\label{eq:logCT}
 \boxed{\displaystyle
 h_{d;\bm m}(a)=\frac{(-1)^{|\bm m|}}{\bm m!}\,
 \CT_{L_N}
 \sum_{N>n_1>\cdots>n_d\ge0}
       \prod_{j=1}^{d}\frac{\log^{m_j}(n_j+a)}{n_j+a}.}
\end{equation}
Here $\CT_{L_N}$ means the constant coefficient of the polynomial
asymptotic, not a bare limit of the divergent sum.
\end{theorem}
\begin{proof}
We give the uniform argument because a formal cutoff comparison alone
would not justify all mixed derivatives. Inductively construct a function
$Q_d(\bm u;L)=K_d(\bm u)+\int_0^L e^{-u_1t}Q_{d-1}(\bm u';t)\dd t$
with the desired approximation. For $d=0$ this is exact. Suppose it is
known at depth $d-1$, with a uniform remainder of size
$O(N^{-1+\sigma}(1+\log N)^M)$ on any compact set with
$\sum|u_j|\le\sigma<1/4$; this estimate is stronger than the claimed
$o(1)$ and closes the induction.

Write $x=N+a$. The increment of
\[
 F_{d,N}-\int_0^{\log(N+a)}e^{-u_1t}Q_{d-1}(\bm u';t)\dd t
\]
is the sum of
\[
 x^{-1-u_1}\bigl(F_{d-1,N}-Q_{d-1}(\bm u';\log x)\bigr)
\]
and the difference between
$x^{-1-u_1}Q_{d-1}(\bm u';\log x)$ and its integral over $[x,x+1]$.
The first expression is $O(x^{-2+\sigma}(1+\log x)^M)$, using the
appropriate tail-variable compact set. Differentiating the second
integrand in $x$ gives the same kind of bound: induction in the defining
integrals gives
$Q_{d-1}=O(x^{\sum_{j\ge2}|u_j|}(1+\log x)^{d-1})$ and the analogous
derivative estimate. Hence these increments are normally summable.
Their partial sums converge to a holomorphic $K_d$, and summing the tail
gives $O(N^{-1+\sigma}(1+\log N)^{M'})$. Cauchy estimates on a slightly
larger compact neighborhood also give convergence of every fixed mixed
derivative.

It remains to identify $K_d$. In the nonempty open domain
$\Re u_j>0$, let $N\to\infty$. Repeatedly integrating
$e^{-u_1t_1}\cdots e^{-u_rt_r}$ over
$0<t_r<\cdots<t_1<\infty$ gives $1/(U_1\cdots U_r)$. Thus the limit
is precisely \eqref{eq:germ} with $K_j$ in place of $\HH_j$.
Inductive uniqueness and the identity theorem imply $K_j=\HH_j$.
This proves \eqref{eq:cutasym} with \eqref{eq:Pcut}.

Finally, Taylor extraction in \eqref{eq:Pcut} replaces each exponential
by finitely many monomials in $t$ and adds one integration at each depth.
This proves the polynomial degree bound and identifies the constant term.
Taylor extraction in the finite sum gives the sign and factorial in
\eqref{eq:logCT}.
\end{proof}

\subsection{The full positive-index logarithmic stuffle algebra}
A letter is a pair $(k,m)$ with $k\ge1$ and $m\ge0$, representing
$x^{-k}\log^m x$. For a word $w=((k_1,m_1),\ldots,(k_d,m_d))$, let
$W_w(N;a)$ be the corresponding finite ordered sum. In the stuffle product,
merging letters means
\begin{equation}\label{eq:merge}
 (k,m)\circ(\ell,n)=(k+\ell,m+n).
\end{equation}

\begin{theorem}[Specified stuffle-preserving regularization]
\label{thm:stuffle}
Every $W_w(N;a)$ has a unique polynomial asymptotic in $L_N$, up to an
error $O(N^{-1}(1+\log N)^M)$ for some $M$. Define $\RR_a(w)$ to be its
constant coefficient. Then
\begin{equation}\label{eq:stuffle}
 \RR_a(w\ast v)=\RR_a(w)\RR_a(v).
\end{equation}
The map agrees with every ordinarily convergent sum in this positive-index
logarithmic class. At depth one,
\begin{equation}\label{eq:letters}
 \RR_a((1,m))=\gamma_m(a),\qquad
 \RR_a((k,m))=(-1)^m\zeta^{(m)}(k,a)\quad(k\ge2).
\end{equation}
For all-one exponents, it is exactly the regularization in
\eqref{eq:logCT}.
\end{theorem}
\begin{proof}
Induct on word length. All finite inner sums grow at most as a power of
$1+\log N$. If the outer exponent $k_1\ge2$, the sum converges and its
tail is $O(N^{-1}(1+\log N)^M)$, so its polynomial asymptotic is constant.
If $k_1=1$, insert the inner polynomial asymptotic into the increment
$\log^{m_1}(N+a)W_{w'}(N;a)/(N+a)$. Integration in $L_N$ gives a
polynomial. The summable remainder and the sum--integral discrepancy are
handled by exactly the increment argument in Theorem \ref{thm:cutoff},
now with no complex exponent perturbations. This proves existence and
the stated error; uniqueness follows because a nonzero polynomial in
$L_N$ cannot tend to zero.

At finite $N$, partitioning the two sets of indices according to their
relative order and coincidences proves
$W_{w\ast v}(N;a)=W_w(N;a)W_v(N;a)$, with merge rule \eqref{eq:merge}.
The remainders still tend to zero after multiplication by any fixed
logarithmic polynomial. The polynomial asymptotics therefore satisfy
the same identity exactly. Evaluation at $L=0$ is multiplicative on
polynomials, proving \eqref{eq:stuffle}. Formula \eqref{eq:letters}
follows from \eqref{eq:g} and the convergent single-index series.
\end{proof}

\begin{remark}[What ``compatible'' does and does not assert]
This prescription is fixed by the harmonic cutoff polynomial and agrees
with the conventional positive-index harmonic regularization at the
corresponding depth-one normalization. It is not a uniqueness theorem
among all conceivable subtraction schemes, nor a claim that raw
multivariate finite parts are path independent. It does not assert an
unmodified shuffle identity for every integral representation. Those
would require separate comparison statements; see \cite{gz,gpxz}.
\end{remark}

\subsection{Every directional coefficient and its counterterm}
Let $c_1,\ldots,c_d\in\mathbb C$ and write
$S_k=c_1+\cdots+c_k$. Assume $S_k\ne0$ for every $k$; individual inner
slopes may vanish.

\begin{corollary}[All straight-ray coefficients]\label{cor:rayall}
For $q\ge0$,
\begin{equation}\label{eq:allray}
 [\eps^q]Z_d(1+c_1\eps,\ldots,1+c_d\eps;a)
 =\sum_{j=1}^{d}
 \frac{[\eps^{q+d-j}]\HH_j(c_{d-j+1}\eps,\ldots,c_d\eps;a)}
      {S_1\cdots S_{d-j}}.
\end{equation}
The leading coefficient is $(S_1\cdots S_d)^{-1}\eps^{-d}$, and the
other negative coefficients follow from the displayed coefficient rule,
with negative Taylor indices interpreted as zero.
\end{corollary}
\begin{proof}
Substitute $u_i=c_i\eps$ in \eqref{eq:germ}. Each prefix denominator supplies $\eps^{d-j}S_1\cdots S_{d-j}$.
\end{proof}

\begin{corollary}[Explicit harmonic counterterm]\label{cor:counterterm}
The counterterm converting the ray constant to the harmonic regularized
all-one value is
\begin{equation}\label{eq:counterterm}
 \boxed{\displaystyle
 \FP Z_d-C_d(a)=
 \sum_{j=1}^{d-1}
 \frac{[\eps^{d-j}]\HH_j(c_{d-j+1}\eps,\ldots,c_d\eps;a)}
      {S_1\cdots S_{d-j}}.}
\end{equation}
Every coefficient on the right is obtained by the finite recursion
\eqref{eq:hrec}. The resulting $C_d(a)$ belongs to the
stuffle-preserving prescription of Theorem \ref{thm:stuffle}.
\end{corollary}

At depth two, this gives
\begin{equation}\label{eq:doubleFP}
 \FP Z_2(1+c\eps,1+d\eps;a)=C_2(a)-\frac dc\gamma_1(a).
\end{equation}
Thus subtracting the ray value alone is not an algebra homomorphism:
the two depth-two orientations carry the extra combined term
$-\gamma_1(a)(d/c+c/d)$. The source's necessary compatibility test is
therefore met exactly, not merely asymptotically \cite{repo-target}.

\section{Ordered depth three and its nonsymmetric coordinate}
\label{sec:triple}

The depth-three specialization of \eqref{eq:germ} is
\begin{equation}\label{eq:triple-germ}
 \boxed{\displaystyle
 Z_3(1+u,1+v,1+w;a)=
 \frac1{u(u+v)(u+v+w)}+
 \frac{g_a(w)}{u(u+v)}+
 \frac{\HH_2(v,w;a)}u+\HH_3(u,v,w;a).}
\end{equation}
The remainder is explicitly
\begin{equation}\label{eq:H3}
 \HH_3(u,v,w;a)=\Atail_3(u,v,w;a)+
 \frac{\HH_2(u+v,w;a)-\HH_2(v,w;a)}u.
\end{equation}
In particular, the lower regular germ is evaluated at the \emph{ordered
suffix} $(v,w)$, not at a symmetrized pair or at $(u+v,w)$ in
\eqref{eq:triple-germ}.

Put
\begin{equation}\label{eq:h10}
 h_{10}=\partial_1\HH_2(0,0;a),\qquad
 h_{01}=\partial_2\HH_2(0,0;a),\qquad
 \Omega(a)=h_{10}-h_{01}.
\end{equation}
The double stuffle identity gives
\begin{equation}\label{eq:Hsym}
 \HH_2(u,v;a)+\HH_2(v,u;a)
 =g_a(u)g_a(v)-\zeta(2+u+v,a),
\end{equation}
so
\begin{equation}\label{eq:hsum}
 h_{10}+h_{01}=-\gamma_0(a)\gamma_1(a)-\zeta'(2,a).
\end{equation}

\begin{theorem}[Every ordered transverse depth-three finite part]
\label{thm:tripleFP}
Let $c\ne0$, $c+d\ne0$, and $c+d+e\ne0$. Then
\begin{align}\label{eq:tripleLaurent}
 Z_3(1+c\eps,1+d\eps,1+e\eps;a)
 ={}&\frac1{c(c+d)(c+d+e)\eps^3}
 +\frac{\gamma_0(a)}{c(c+d)\eps^2}\notag\\
 &+\frac1\eps\left\{\frac{C_2(a)}c-
          \frac{e\gamma_1(a)}{c(c+d)}\right\}
 +F_{c,d,e}(a)+O(\eps),
\end{align}
where
\begin{align}\label{eq:tripleFP}
 \boxed{\displaystyle F_{c,d,e}(a)
 = C_3(a)+\frac{d h_{10}+e h_{01}}c
              +\frac{e^2\gamma_2(a)}{2c(c+d)}.}
\end{align}
Equivalently,
\begin{align}\label{eq:tripleOmega}
 F_{c,d,e}(a)={}&C_3(a)
 -\frac{d+e}{2c}\bigl(\gamma_0(a)\gamma_1(a)+\zeta'(2,a)\bigr)\notag\\
 &+\frac{d-e}{2c}\Omega(a)
 +\frac{e^2\gamma_2(a)}{2c(c+d)}.
\end{align}
\end{theorem}
\begin{proof}
Expand $g_a(e\eps)$ through degree two and
$\HH_2(d\eps,e\eps;a)$ through degree one in
\eqref{eq:triple-germ}, then use \eqref{eq:Cdef} and
\eqref{eq:hsum}. Holomorphy of the remainders controls the omitted terms.
\end{proof}

\paragraph{An exact reduced two-parameter family.}
When the two inner slopes agree, the nonsymmetric coordinate disappears:
\begin{equation}\label{eq:equalinner}
 F_{c,d,d}(a)=C_3(a)-\frac dc
      \bigl(\gamma_0(a)\gamma_1(a)+\zeta'(2,a)\bigr)
       +\frac{d^2\gamma_2(a)}{2c(c+d)}.
\end{equation}
In particular,
\begin{equation}\label{eq:diagonaltriple}
 F_{1,1,1}(a)=C_3(a)-\gamma_0(a)\gamma_1(a)
                    -\zeta'(2,a)+\frac{\gamma_2(a)}4.
\end{equation}
This also follows independently by expanding
\[
 Z_3(s,s,s;a)=\frac{\zeta(s,a)^3-3\zeta(s,a)\zeta(2s,a)
                                  +2\zeta(3s,a)}6.
\]

\subsection{A convergent Stieltjes formula for the orientation constant}

\begin{proposition}[Exact nonsymmetric coordinate]\label{prop:omega}
For every $a>0$,
\begin{equation}\label{eq:omega-series}
 \boxed{\displaystyle
 \Omega(a)= -\frac{\gamma_2(a)}2+
 \sum_{n=0}^\infty
 \frac{\log(n+a)\gamma_0(n+a)-\gamma_1(n+a)
                    +\tfrac12\log^2(n+a)}{n+a}.}
\end{equation}
The series is ordinarily absolutely convergent. Its summand is
$1/(12x^3)+O(x^{-5})$ as $x=n+a\to\infty$.
\end{proposition}
\begin{proof}
From \eqref{eq:Hrec} at depth two,
\[
 h_{10}=\partial_1\Atail_2(0,0;a)+g_2(a),\qquad
 h_{01}=\partial_2\Atail_2(0,0;a)+2g_2(a).
\]
For $x=n+a$, the difference of the tail derivatives is
\[
 \frac{B_u(0,x)+\log x\,B(0,x)}x.
\]
The Laurent expansion of \eqref{eq:B} gives
$B(0,x)=\gamma_0(x+1)+\log x$ and
$B_u(0,x)=-\gamma_1(x+1)-\tfrac12\log^2x$.
The shifts
$\gamma_0(x+1)=\gamma_0(x)-1/x$ and
$\gamma_1(x+1)=\gamma_1(x)-\log x/x$ cancel the extra terms, proving
\eqref{eq:omega-series}.

For the printed leading term, scale $t=y/x$ in \eqref{eq:Bint}
before differentiating. The globally bounded Bernoulli remainder for
$K(y/x)$ gives, locally uniformly with all fixed $u$-derivatives,
\[
 x^uB(u,x)=-\frac1{2x}+\frac{1+u}{12x^2}
       -\frac{(1+u)(2+u)(3+u)}{720x^4}+O(x^{-6}).
\]
For example, the remainder after the cubic term is bounded by a constant
times $(y/x)^5$ for all $y>0$, and the $e^{-y}$ factor is integrable
against every required logarithmic power. Differentiating the scaled
expression at zero gives
$B_u(0,x)+\log x\,B(0,x)=x^{-2}/12-11x^{-4}/720+O(x^{-6})$.
Division by $x$ proves the asserted convergence and leading term.
Normal differentiation in Lemma \ref{lem:normal} justifies the series
operations as well.
\end{proof}

This solves the requested construction of a normalized nonsymmetric
coordinate, but not its reduction to any smaller prescribed constant
field. In the notation of the source's convergent sums,
$\Omega=S_{1,0}-S_{0,1}-\gamma_2(a)/2$. It is essential to keep the
last Stieltjes term: dropping it changes the actual ordered finite part.

\subsection{The full labelled symmetrization check}
Let $c_1=c,c_2=d,c_3=e$ and assume all three slopes and all pair and triple
sums are nonzero. Summing over all six labelled permutations yields
\begin{align}\label{eq:symcheck}
 \sum_{\sigma\in S_3}F_{c_{\sigma(1)},c_{\sigma(2)},c_{\sigma(3)}}(a)
 ={}&6C_3(a)
 -\bigl(\gamma_0\gamma_1+\zeta'(2,a)\bigr)
          \sum_{i\ne j}\frac{c_i}{c_j}\notag\\
 &+\frac{\gamma_2(a)}2
       \left(\frac{c^2}{de}+\frac{d^2}{ce}+\frac{e^2}{cd}\right).
\end{align}
The $\Omega$ terms cancel exactly. To verify independently, expand the
finite-stuffle continuation
\begin{align*}
 &\sum_{\sigma\in S_3}Z_3(s_{\sigma(1)},s_{\sigma(2)},s_{\sigma(3)};a)\\
 &\quad=\prod_i\zeta(s_i,a)
 -\sum_{\{i,j,k\}=\{1,2,3\}\!,\,i<j}\zeta(s_i+s_j,a)\zeta(s_k,a)
 +2\zeta(s_1+s_2+s_3,a).
\end{align*}
Each pair is counted once in the middle sum. This check explains why a
fully symmetrized computation cannot by itself determine the ordered
constant $\Omega$.

\section{Nonlinear and tangential paths: exactly how much data is needed}
\label{sec:curves}

Let $u_j(\eps)$ be analytic with $u_j(0)=0$, and let
$U_k(\eps)=\sum_{j\le k}u_j(\eps)$. Suppose no $U_k$ is identically
zero. Write
\begin{equation}\label{eq:valuations}
 \nu_k=\ord_0U_k\ge1,\qquad
 P_k=\nu_1+\cdots+\nu_k,\qquad P=P_d.
\end{equation}

\begin{theorem}[Finite-jet algorithm for every admissible analytic path]
\label{thm:curve}
The substituted germ has pole order at most $P$. Its finite part is
obtained by substituting the curve into \eqref{eq:germ} and performing
finite power-series arithmetic. It is sufficient to know each prefix
$U_k$ through degree $\nu_k+P$, together with finitely many regular
coefficients provided by \eqref{eq:hrec}. The contribution with denominator
$U_1\cdots U_r$ needs only that denominator's normalized reciprocal through
degree $P_r$ and its numerator through degree $P_r$ after substitution.
For transverse curves, $\nu_k=1$ for all $k$, and the $(d+1)$-jet of the
curve suffices. The degree requirement $\nu_k+P$ is sharp in general for
the prefix data.
\end{theorem}
\begin{proof}
Factor $U_k(\eps)=\eps^{\nu_k}V_k(\eps)$ with $V_k(0)\ne0$.
The reciprocal of a prefix product is
$\eps^{-P_r}\prod_{k\le r}V_k(\eps)^{-1}$. Its constant contribution
against a holomorphic numerator requires only terms through degree $P_r$
in the remaining product. Reciprocal-series recursion uses $V_k$ through
that degree, equivalently $U_k$ through degree $\nu_k+P_r$. Taking
$r\le d$ proves sufficiency. The Taylor series of every numerator is
holomorphic, so only finitely many of its coefficients can enter.

For sharpness, take all $V_i$ constant except one and change
$V_k$ by $b\eps^P$. In the leading term $1/(U_1\cdots U_d)$ this changes
the finite part by a nonzero multiple of $b$. Every other term has a
strictly smaller denominator valuation and cannot detect this high-degree
change. The $U_k$ are independent coordinates, so these choices correspond
to actual curves $u_j=U_j-U_{j-1}$.
\end{proof}

This theorem includes tangencies of arbitrarily high finite order.
Curves lying \emph{identically} in a polar divisor remain excluded:
restriction of a meromorphic germ there is not defined without an
additional operation. The theorem is not a prescription for that different
problem.

\subsection{Two explicit depth-three consequences}
For a straight ray, write the coefficients in \eqref{eq:tripleLaurent} as
$A_{-3},A_{-2},A_{-1},F$. Reparametrize the same geometric ray by
\[
 u_j=c_j f(\eps),\qquad
 f(\eps)=\eps+\alpha\eps^2+\beta\eps^3+\chi\eps^4+O(\eps^5).
\]
Then
\begin{equation}\label{eq:reparam}
 \FP Z_3=F-\alpha A_{-1}+(3\alpha^2-2\beta)A_{-2}
              +(-10\alpha^3+12\alpha\beta-3\chi)A_{-3}.
\end{equation}
This follows by expanding $f^{-1},f^{-2},f^{-3}$ and exhibits why the
fourth jet is genuinely needed.

A tangential example has
\[
 u=\eps,\qquad v=-\eps+\eps^2,\qquad w=\eps.
\]
Its prefix valuations are $(1,2,1)$, so the pole order is four. Direct
substitution in \eqref{eq:triple-germ} gives the exact constant
\begin{equation}\label{eq:tangent}
 \boxed{\displaystyle
 \FP Z_3(1+\eps,1-\eps+\eps^2,1+\eps;a)
   =C_3(a)-\Omega(a)-\frac{\gamma_3(a)}6+1.}
\end{equation}
Indeed the leading term is $\eps^{-4}/(1+\eps)$, the second contributes
$g_3(a)$, and $\HH_2(-\eps+\eps^2,\eps)/\eps$ contributes
$-h_{10}+h_{01}$. The additional rational constant is a real curve effect,
not an arbitrary regularization constant.

\section{Polylogarithmic Mellin coordinates for every regular jet}
\label{sec:mellin}

For $0<z<1$, define the common-shift multiple polylogarithm
\begin{equation}\label{eq:Lshift}
 \LL_{s_1,\ldots,s_r}(z;a)
 =\sum_{n_1>\cdots>n_r\ge0}z^{n_1+a}
                       \prod_{j=1}^{r}(n_j+a)^{-s_j}.
\end{equation}
The exponential outer weight makes this entire in the order variables.
At $a=1$ it is the standard multiple polylogarithm
$\Li_{s_1,\ldots,s_r}(z,1,\ldots,1)$, with the outer index largest.
For a single index it is $z^a\Phi(z,s,a)$; at $a=1$ this
reduces to the classical polylogarithm convention of \cite{dlmf-polylog}.

\begin{theorem}[Ordinary Mellin representation]\label{thm:Mellin}
Near $\bm u=0$, every tail term has the absolutely convergent integral
\begin{equation}\label{eq:AMellin}
 \boxed{\displaystyle
 \Atail_d(\bm u;a)=\frac1{\Gamma(1+u_1)}
 \int_0^\infty t^{u_1}K(t)
          \LL_{1+u_2,\ldots,1+u_d}(e^{-t};a)\dd t.}
\end{equation}
Every fixed mixed order derivative may be taken under the integral.
\end{theorem}
\begin{proof}
Insert \eqref{eq:Bint} into \eqref{eq:Ad}. For real parameters in a small
neighborhood of zero, absolute convergence follows either from the
summable bounds in Lemma \ref{lem:normal} applied to the absolute integral,
or directly by summing the exponentially weighted inner finite harmonic
sums. Near $t=0$, a bound of the form
$C t^{-\sigma}(1+|\log t|)^{d-1}$ suffices for the polylogarithm, while
$K$ is bounded; choose the neighborhood with total exponent loss below
$1/2$. At infinity the factor $e^{-at}$ dominates every polynomial.
Fixed derivatives add only logarithmic powers and obey the same argument.
Fubini and holomorphic continuation prove the assertion.
\end{proof}

Define the explicit polynomials $P_m(X)$ by
\begin{equation}\label{eq:Pm}
 \frac{t^u}{\Gamma(1+u)}=\sum_{m\ge0}P_m(\log t)u^m,
\end{equation}
so
\begin{align}
 P_0(X)&=1,\quad P_1(X)=X+\gamma,\notag\\
 P_2(X)&=\frac{(X+\gamma)^2-\zeta(2)}2,\label{eq:Psmall}\\
 P_3(X)&=\frac{(X+\gamma)^3-3\zeta(2)(X+\gamma)+2\zeta(3)}6.
 \notag
\end{align}
The $\gamma$ in these universal polynomials is Euler's constant, not
$\gamma_0(a)$ when $a\ne1$.

\begin{corollary}[Explicit coefficient integrals]\label{cor:Acoeff}
For every nonnegative multi-index,
\begin{align}\label{eq:Acoeff}
 A_{d;\bm m}(a)=\frac1{\prod_{j=2}^{d}m_j!}
 \int_0^\infty K(t)P_{m_1}(\log t)
 \left.
 \partial_{s_2}^{m_2}\cdots\partial_{s_d}^{m_d}
 \LL_{s_2,\ldots,s_d}(e^{-t};a)
 \right|_{s_2=\cdots=s_d=1}\dd t.
\end{align}
Together with \eqref{eq:hrec}, these are finite integral coordinates for
every ordered regular Hurwitz jet.
\end{corollary}

\subsection{Rational shifts have explicit cyclotomic coordinates}
Use the independent-color convention
\[
 \Li_{\bm s}(z_1,\ldots,z_r)
   =\sum_{\ell_1>\cdots>\ell_r\ge1}
           \prod_{j=1}^{r}\frac{z_j^{\ell_j}}{\ell_j^{s_j}}.
\]
If $a=p/q$, $1\le p\le q$, and $\omega=e^{2\pi i/q}$, then
\begin{equation}\label{eq:rationalfilter}
 \boxed{\displaystyle
 \LL_{\bm s}(z;p/q)
 =q^{s_1+\cdots+s_r-r}
  \sum_{j_1,\ldots,j_r=0}^{q-1}\omega^{-p(j_1+\cdots+j_r)}
  \Li_{\bm s}(\omega^{j_1}z^{1/q},\omega^{j_2},\ldots,\omega^{j_r}).}
\end{equation}
Indeed, put $\ell_i=qn_i+p$ and insert the elementary residue-class
indicator in every index. Strict ordering is preserved. The series are
absolutely convergent, and every order derivative is obtained by
ordinary differentiation, including the powers of $\log q$ from the
prefactor. The shift recurrence in Section \ref{sec:shifts} handles larger
positive rational shifts. This is an exact alphabet statement, not a
minimality or period-independence assertion.

\section{An all-depth Mellin hierarchy and two closed generating identities}
\label{sec:master}

The leading-axis specialization of \eqref{eq:hrec} telescopes particularly
well. Repeatedly take $m_2=\cdots=m_d=0$; every binomial coefficient is
one. The result is
\begin{equation}\label{eq:axis}
 C_d(a)=g_{d-1}(a)+\sum_{r=2}^{d}A_{r;(d-r,0,\ldots,0)}(a).
\end{equation}

\begin{theorem}[All-depth exact Mellin identities]\label{thm:hierarchy}
For every integer $d\ge2$ and every $a>0$,
\begin{equation}\label{eq:hierarchy}
 \boxed{\displaystyle
 \int_0^\infty K(t)\sum_{r=2}^{d}P_{d-r}(\log t)
          \LL_{\{1\}^{r-1}}(e^{-t};a)\dd t
   =C_d(a)-\frac{(-1)^{d-1}\gamma_{d-1}(a)}{(d-1)!}.}
\end{equation}
All the integrals are ordinary absolutely convergent integrals.
\end{theorem}
\begin{proof}
Combine \eqref{eq:axis} with Corollary \ref{cor:Acoeff}. The sum over
$r$ is finite, so no additional interchange is needed.
\end{proof}

For $a=1$, let $L(t)=-\log(1-e^{-t})$. The elementary all-one identity
$\Li_{\{1\}^{m}}(z,1,\ldots,1)=(-\log(1-z))^m/m!$ follows by
successive differentiation in $z$ and the zero boundary value. Hence
\eqref{eq:hierarchy} becomes
\begin{equation}\label{eq:elementaryhierarchy}
 \int_0^\infty K(t)\sum_{r=2}^{d}
    P_{d-r}(\log t)\frac{L(t)^{r-1}}{(r-1)!}\dd t
 =C_d(1)-\frac{(-1)^{d-1}\gamma_{d-1}}{(d-1)!}.
\end{equation}
The first three identities are
\begin{align}
 \int_0^\infty K(t)L(t)\dd t
     &=\frac{\gamma^2-\zeta(2)}2+\gamma_1,\label{eq:I2}\\
 \int_0^\infty K(t)\left((\log t+\gamma)L(t)+\frac{L(t)^2}2\right)\dd t
     &=\frac{\gamma^3-3\gamma\zeta(2)+2\zeta(3)}6-\frac{\gamma_2}2,
       \label{eq:I3}\\
 \int_0^\infty K(t)\left(P_2(\log t)L(t)
        +(\log t+\gamma)\frac{L(t)^2}2+\frac{L(t)^3}6\right)\dd t
     &=C_4(1)+\frac{\gamma_3}6.\label{eq:I4}
\end{align}
These are explicit identities rather than inequalities or asymptotic
approximations.

\begin{theorem}[Elementary master identity]\label{thm:master}
For every $z\in\mathbb C$ with $\Re z>-1$,
\begin{equation}\label{eq:master}
 \boxed{\displaystyle
 \int_0^\infty K(t)t^z\bigl((1-e^{-t})^{-z}-1\bigr)\dd t
       =\frac1z-\Gamma(1+z)\zeta(1+z),}
\end{equation}
with the removable value zero at $z=0$. Every fixed $z$-derivative is
obtained by differentiation under the integral on this half-plane.
\end{theorem}
\begin{proof}
The difference in the integrand is essential. At zero it is bounded by
$C(1+t^{\Re z})$ on compact $z$-sets, and at infinity it is
$O(t^{\Re z}e^{-t})$ before multiplication by $K$. These bounds and their
logarithmic derivatives prove holomorphy for $\Re z>-1$.

For $\Re z>0$, put $Q_z(t)=(t/(1-e^{-t}))^z$. Then
$Q_z'(t)=-zK(t)Q_z(t)$, and direct rearrangement gives
\[
 K(t)(Q_z(t)-t^z)
 =-\frac1z\frac{d}{dt}(Q_z(t)-t^z)-\frac{t^z}{e^t-1}.
\]
At infinity $Q_z(t)-t^z\to0$; at zero it tends to one. Integration
therefore gives $1/z-\Gamma(1+z)\zeta(1+z)$ by the ordinary zeta Mellin
integral. Analytic continuation proves the full stated half-plane. At
zero the residue and constant terms of $\Gamma(1+z)\zeta(1+z)$ are
$1/z$ and zero, respectively, proving the removable value. The local
integrable bounds justify all order derivatives.
\end{proof}

Multiplying \eqref{eq:master} by $z/\Gamma(1+z)$ and extracting
$z^d$ proves \eqref{eq:elementaryhierarchy} independently of the
multiple-zeta germ. Thus the hierarchy has both a nested-sum proof and
a one-variable calculus proof. This independent derivation is useful for
checking signs and Stieltjes factorials.

\subsection{A common-Hurwitz-shift hypergeometric master}
Define, initially for sufficiently small $|z|$,
\begin{equation}\label{eq:Gat}
 G_a(z,t)=\sum_{m\ge1}z^m\LL_{\{1\}^{m}}(e^{-t};a).
\end{equation}
Finite elementary symmetric sums give an explicit closed form:
\begin{equation}\label{eq:Ghyper}
 G_a(z,t)=\frac za e^{-at}
     {}_2F_1(1,a+z;a+1;e^{-t}).
\end{equation}
To check the coefficient, fix the largest index $n$. Its contribution is
\[
 \frac{z e^{-(n+a)t}}{n+a}
     \prod_{j=0}^{n-1}\left(1+\frac z{j+a}\right)
 =\frac za e^{-(n+a)t}\frac{(a+z)_n}{(a+1)_n}.
\]
Summing over $n$ gives \eqref{eq:Ghyper}.

\begin{corollary}[Hurwitz-shifted master identity]\label{cor:hyper}
For $a>0$ and $\Re z>-1$,
\begin{equation}\label{eq:hypermaster}
 \boxed{\displaystyle
 \frac z{\Gamma(1+z)}\int_0^\infty K(t)t^z G_a(z,t)\dd t
      =\frac{\Gamma(a)}{\Gamma(a+z)}-1-zg_a(z).}
\end{equation}
The value at $z=0$ is removable. Every fixed $z$-derivative may be
taken under the integral. At $a=1$, this is exactly \eqref{eq:master}
after cancellation of the prefactor. Its Taylor coefficients are
\eqref{eq:hierarchy}.
\end{corollary}
\begin{proof}
We give an independent integral proof that also establishes the full
domain. Put $q=e^{-t}$ and use the incomplete beta integral
\[
 B_q(a,z)=\int_0^q x^{a-1}(1-x)^{z-1}\dd x.
\]
For $q<1$, this is entire in $z$. Euler's hypergeometric transformation,
or direct comparison of power-series coefficients, gives
\begin{equation}\label{eq:Gbeta}
 G_a(z,t)=z(1-q)^{-z}B_q(a,z).
\end{equation}
Differentiating this ordinary integral and its prefactor yields
\[
 \partial_tG_a(z,t)=-\frac{z}{e^t-1}G_a(z,t)
                        -\frac{z e^{-at}}{1-e^{-t}}.
\]
Therefore, with $F(t)=t^zG_a(z,t)$,
\begin{equation}\label{eq:hyperderivative}
 K(t)F(t)=-\frac{F'(t)}z-\frac{t^ze^{-at}}{1-e^{-t}}.
\end{equation}
For $\Re z>0$, the boundary values are $F(\infty)=0$ and
\[
 F(0)=zB(a,z)=\frac{\Gamma(a)\Gamma(1+z)}{\Gamma(a+z)}.
\]
Integrating \eqref{eq:hyperderivative} and using the Hurwitz Mellin
integral proves \eqref{eq:hypermaster} on this open half-plane.

To extend the ordinary integral, set $w=1-e^{-t}$ in \eqref{eq:Gbeta}:
\[
 F(t)=z\left(\frac{t}{w}\right)^z
       \int_w^1 (1-y)^{a-1}y^{z-1}\dd y.
\]
On each compact $z$-set with $\Re z>-1$, this is bounded near zero by
$C(1+t^{-\rho})(1+|\log t|)$ for some $\rho<1$. The same estimate
with extra logarithmic powers holds for fixed $z$-derivatives. At
infinity, $G_a(z,t)=O(e^{-at})$ locally uniformly. Because $K$ is bounded
at zero, these bounds prove holomorphy of the integral and differentiation
under it throughout $\Re z>-1$. The right side is holomorphic there:
$\Gamma(a)/\Gamma(a+z)$ is entire in $z$, and $zg_a(z)$ is entire.
The identity theorem proves the stated extension, including its value
at zero. Finally, coefficient extraction in
$z t^zG_a(z,t)/\Gamma(1+z)$ gives \eqref{eq:hierarchy}. This also
independently verifies the nested-sum derivation for every $a>0$.
\end{proof}

\subsection{Two additional mixed order-derivative identities}
The first-order coefficients of \eqref{eq:Hsym}, combined with
\eqref{eq:Acoeff}, give
\begin{align}\label{eq:mixedMellin}
 &\int_0^\infty K(t)\left\{(\log t+\gamma)\LL_1(e^{-t};a)
       +\left.\partial_s\LL_s(e^{-t};a)\right|_{s=1}\right\}\dd t
 \notag\\[-2pt]
 &\hspace{25mm}=-\gamma_0(a)\gamma_1(a)-\zeta'(2,a)
                                      -\frac32\gamma_2(a),
\end{align}
and the nonsymmetric coordinate has the ordinary integral
\begin{align}\label{eq:omegaMellin}
 \Omega(a)={}&-\frac{\gamma_2(a)}2+
 \int_0^\infty K(t)\left\{(\log t+\gamma)\LL_1(e^{-t};a)
       -\left.\partial_s\LL_s(e^{-t};a)\right|_{s=1}\right\}\dd t.
\end{align}
At $a=1$, no order derivative is needed for a second representation.
Using \eqref{eq:I3} and \eqref{eq:hsum},
\begin{equation}\label{eq:omegaelementary}
 \boxed{\displaystyle
 \Omega(1)=2C_3(1)+\gamma\gamma_1+\zeta'(2)
                         -\int_0^\infty K(t)L(t)^2\dd t.}
\end{equation}
Equations \eqref{eq:omega-series}, \eqref{eq:omegaMellin}, and
\eqref{eq:omegaelementary} are exact coordinate changes. None asserts that
$\Omega$ is independent of, or reducible to, a specified smaller field.

\section{Shift transport, harmonic recurrences, and polygamma primitives}
\label{sec:shifts}

\begin{proposition}[Exact shift recurrence of every regular germ]
\label{prop:shift}
For $d\ge1$,
\begin{equation}\label{eq:Hshift}
 \boxed{\displaystyle
 \HH_d(u_1,\ldots,u_d;a)-\HH_d(u_1,\ldots,u_d;a+1)
  =a^{-1-u_d}\HH_{d-1}(u_1,\ldots,u_{d-1};a+1).}
\end{equation}
In coefficients,
\begin{equation}\label{eq:hshift}
 h_{d;\bm m}(a)-h_{d;\bm m}(a+1)
 =\frac{(-\log a)^{m_d}}{a\,m_d!}
          h_{d-1;(m_1,\ldots,m_{d-1})}(a+1).
\end{equation}
\end{proposition}
\begin{proof}
In the convergence domain, separating the terms with smallest index zero
gives
\[
 Z_d(\bm s;a)-Z_d(\bm s;a+1)
          =a^{-s_d}Z_{d-1}(s_1,\ldots,s_{d-1};a+1).
\]
Insert \eqref{eq:germ}. Inductively apply \eqref{eq:Hshift} to every
lower-depth suffix numerator. Their prefix-denominator terms match
those on the right exactly; the only unmatched holomorphic term is the
one displayed in \eqref{eq:Hshift}. The depth-one base is the usual
Hurwitz shift. Extracting the coefficient of $a^{-u_d}$ proves
\eqref{eq:hshift}.
\end{proof}

Repeated use of \eqref{eq:hshift} gives finite harmonic identities at
integer shifts. For instance,
\begin{equation}\label{eq:Cshift}
 C_d(a)-C_d(a+1)=\frac1a C_{d-1}(a+1),
\end{equation}
which also follows from the Gamma functional equation. The nonsymmetric
coordinate obeys
\begin{equation}\label{eq:Omegashift}
 \Omega(a)-\Omega(a+1)
   =\frac{-\gamma_1(a+1)+\log a\,\gamma_0(a+1)}a.
\end{equation}
This is an exact discrete antiderivative relation for a nonlinear
Stieltjes expression.

\begin{theorem}[All-depth polygamma primitives]\label{thm:primitive}
For $d\ge1$,
\begin{equation}\label{eq:Cderivative}
 C_d'(a)=\sum_{k=1}^{d}(-1)^k\zeta(k+1,a)C_{d-k}(a).
\end{equation}
Consequently, for $0<a<b$,
\begin{equation}\label{eq:primitive}
 \boxed{\displaystyle
 \int_a^b\sum_{k=1}^{d}(-1)^k\zeta(k+1,x)C_{d-k}(x)\dd x
                       =C_d(b)-C_d(a).}
\end{equation}
Every integrand and endpoint is a finite polygamma polynomial.
\end{theorem}
\begin{proof}
Differentiate \eqref{eq:gamma} in $a$. The logarithmic derivative is
\[
 \psi(a)-\psi(a+z)
 =-\sum_{k\ge1}\frac{\psi^{(k)}(a)}{k!}z^k
 =\sum_{k\ge1}(-1)^k\zeta(k+1,a)z^k.
\]
Coefficient comparison gives \eqref{eq:Cderivative}; ordinary integration
on $[a,b]$ gives \eqref{eq:primitive}. All functions are smooth there.
\end{proof}

For example, an explicit cubic primitive is
\begin{align}\label{eq:cubicprimitive}
 &\int\left(-\zeta(2,a)\frac{\gamma_0(a)^2-\zeta(2,a)}2
              +\gamma_0(a)\zeta(3,a)-\zeta(4,a)\right)\dd a\notag\\
 &\qquad=\frac{\gamma_0(a)^3-3\gamma_0(a)\zeta(2,a)+2\zeta(3,a)}6
          +\text{constant}.
\end{align}
The point is the uniform all-depth generator, not a claim that the
individual low-degree Gamma differentiation rules are new.

\section{Audit findings and reproducible verification}
\label{sec:audit}

\subsection{Corrections and scope safeguards}
No substantive error was found in the specific current repository
proof sections examined for this continuation. In particular, the
source's double germ and its directional formula agree with
\eqref{eq:germ} and \eqref{eq:doubleFP}. The following are concrete
integration safeguards, not fabricated errata.

\begin{enumerate}[itemsep=5pt]
\item \textbf{Keep ordered suffixes.} The numerator of the simple outer
pole in \eqref{eq:triple-germ} is $\HH_2(v,w)$.
Replacing it by its symmetric average loses $\Omega$ and changes ordered
ray constants. Replacing it by $\HH_2(u+v,w)$ requires the compensating
holomorphic divided difference in \eqref{eq:H3}.
\item \textbf{Do not identify raw finite parts with harmonic constants.}
Already at depth two the difference is $-(d/c)\gamma_1(a)$. Equations
\eqref{eq:counterterm} and \eqref{eq:tripleOmega} give the required
higher-depth corrections. A constant-term operation on Laurent series
is not automatically multiplicative.
\item \textbf{Do not specify only a tangent direction for a nonlinear
curve.} Depth three generally reads the fourth jet, and tangencies can
raise the required order further. Curves contained in a polar divisor
need a separate restriction prescription.
\item \textbf{Translate conventions before comparing papers.}
The ordering in \cite{mow} is reversed relative to \eqref{eq:Zd}, and
its one-variable Laurent coefficients are our
$(-1)^m\gamma_m/m!$. Its analytic remainders should not be identified
with the harmonic-compatible $\HH_d$ without the corresponding divided
differences.
\item \textbf{Preserve the subtraction inside the Mellin integral.}
The two pieces of $K(t)$ in \eqref{eq:I2} or the two terms in the
bracket of \eqref{eq:master} need not have the same convergence domain
separately. The displayed grouped integrals are the proved identities.
\item \textbf{Retain arithmetic status.} None of the results proves
$S_6$, the current revised $S_8$, an arithmetic minimality claim for
$\Omega$, or independence of a selected period basket. No such status
should be promoted during integration.
\end{enumerate}

\subsection{Independent checks}
The packaged program \src{verification/verify.py} has separate exact,
integral, and ray suites. The exact suite checks divided-difference
coefficients, the complete six-permutation depth-three expression,
nonlinear reparametrization, a tangential constant, Gamma coefficient
normalizations, finite stuffles with arbitrary rational logarithm labels,
and root-of-unity residue indicators. A deliberately corrupted identity
is rejected.

The integral suite checks the master identity at real and complex orders,
the common-shift hypergeometric master, the elementary hierarchy through
depth six, the orientation constant by an independent digamma-tail sum,
Gamma coefficients by Cauchy extraction, and polygamma primitive
identities. The ray suite computes two nonsymmetric triple finite parts
by Cauchy extraction from a separately implemented Euler--Maclaurin
nested-zeta evaluator; it does not use the $\HH_d$ polar formula to
evaluate the function under test. A separate cutoff comparison tests
that evaluator away from the pole. The evaluator and its evidential
limitations are described in Appendix \ref{app:numerical}.

At $a=1$, the numerical coordinates are
\begin{align}
 h_{10}&\simeq0.536787024358418567503607884348811674,\notag\\
 h_{01}&\simeq0.442791676623652925577095734313457621,\label{eq:numbers}\\
 \Omega(1)&\simeq0.093995347734765641926512150035354053.\notag
\end{align}
These decimal values are diagnostics, not certified enclosures. Exact
execution counts, tolerances, actual residuals, software versions, and
run times are recorded in \src{verification/results_*.json}. The
accompanying verification summary reports the completed runs, rather
than counting unexecuted test plans.

The completed default run passes 278 exact assertions, rejects one
deliberate corruption, and passes 33 floating-point comparisons at
55 working decimal digits. The 30 integral/Gamma comparisons have maximum
observed absolute discrepancy below $3\cdot10^{-38}$; the three
independent ray/cutoff comparisons have maximum discrepancy below
$1.2\cdot10^{-32}$. These are observed numerical residuals, not proved
error bounds. The record includes negative and complex orders and the
exceptional shifted parameter $a+z=0$.

\subsection{Reproduction and integration}
The article builds without any repository files:
\begin{verbatim}
latexmk -pdf -interaction=nonstopmode -halt-on-error article.tex
python verification/verify.py --part all
\end{verbatim}
The suggested insertion is a continuation of the ordered Hurwitz-germ
material, with the Mellin hierarchy cross-referenced from the
polylogarithm and Stieltjes chapters. A short theorem fragment and precise
status notes are in \src{integration/}. No remote files were modified.

\section{Further research questions}
\label{sec:future}

\begin{enumerate}[label=\textbf{Q\arabic*.},itemsep=9pt]
\item \textbf{Arithmetic coordinates for the first orientation constant.}
For $a=1$, $1/2$, and $1/3$, specify a finite comparison algebra of
Gamma, zeta-derivative, Dirichlet-$L$-derivative, and cyclotomic
polylogarithm values. Determine whether \eqref{eq:omega-series} or
\eqref{eq:omegaelementary} reduces to it. The new ordinary integral is
an exact coordinate, not an irreducibility proof.

\item \textbf{Representation-theoretic organization of higher regular
jets.} The first depth-three directional asymmetry is one-dimensional.
At depth four, classify the permutation components of the lower-depth
jets entering \eqref{eq:counterterm}; identify which components are
forced by stuffle identities and which require genuinely ordered
coordinates. The finite coefficient algorithm makes this a concrete
linear-algebra problem at each jet order.

\item \textbf{A complete comparison with other renormalization schemes.}
Construct an explicit character transformation from the harmonic cutoff
map in Theorem \ref{thm:stuffle} to a specified Guo--Zhang or
Manchon--Paycha prescription. Match regulator, depth-one normalization,
ordering, and logarithmic letters before comparing. This paper supplies
the raw directional-to-harmonic counterterms, not that full comparison.

\item \textbf{Curves contained in polar divisors.} Theorem
\ref{thm:curve} handles every finite-order tangency off the divisors.
For a path contained in $U_k=0$, compare iterated residue removal,
Hadamard subtraction in a chosen transverse coordinate, and a
renormalized restriction. Determine precise compatibility conditions;
there is no unqualified meromorphic restriction to start from.

\item \textbf{Unequal Hurwitz shifts and colored logarithmic letters.}
Extend \eqref{eq:Hrec} to distinct shifts $a_1,\ldots,a_d$ and
independent root-of-unity colors. The common-shift outer-tail collapse
then has a nontrivial two-shift correction. Seek an explicit ordinary
integral recursion that retains a useful cutoff-algebra interpretation.

\item \textbf{Global continuation of the shifted master identity.}
The full half-plane $\Re z>-1$ is established in
\eqref{eq:hypermaster}. Continue the integral past the negative-integer
endpoint singularities using explicit local subtractions, and classify
exceptional parameter values where a residue vanishes. Match these
subtractions to higher negative-order Hurwitz and Stieltjes derivatives.
The elementary $a=1$ identity gives an exact model for this analysis.

\item \textbf{Classical reductions of rational-shift Mellin coefficients.}
At small denominator $q$, use \eqref{eq:rationalfilter} to reduce low-order
instances of \eqref{eq:Acoeff} to a fixed cyclotomic iterated-integral
alphabet. Track order derivatives explicitly instead of differentiating
an integer-index finite reduction outside its proved range.

\item \textbf{Formal and certified implementations.} Formalize the
normal-convergence lemma, divided-difference recursion, harmonic cutoff
comparison, and finite stuffle algebra separately. Replace the numerical
Euler--Maclaurin and quadrature diagnostics by rigorous remainder
intervals. The existing corruption control tests the exact verifier,
not the analytic proof.

\item \textbf{The Gaussian arithmetic targets.} Pursue a separate exact
functional certificate for $S_6$, then for the current $S_8$ candidate.
The present local germ calculus can organize constants and regulators,
but no argument here makes it supply the missing arithmetic relation.
The source's conjectural labels remain appropriate.
\end{enumerate}

\appendix
\section{Details of the independent numerical ray evaluator}
\label{app:numerical}

The numerical diagnostics use a cutoff $A=N+a$. The finite splitting
identity for depth three is exact:
\begin{align}\label{eq:split}
 Z_3(s,t,v;a)={}&F_{3,N}(s,t,v;a)
 +\zeta(s,A)F_{2,N}(t,v;a)\notag\\
 &+Z_2(s,t;A)F_{1,N}(v;a)+Z_3(s,t,v;A).
\end{align}
In this appendix only, the finite sums are written with their actual
exponents rather than the perturbations in \eqref{eq:FN}. The identity
partitions an ordered tuple by the number of initial indices at least
$N$.

The large-$A$ tails are evaluated by the independent Euler--Maclaurin
truncations
\begin{align}\label{eq:EM2}
 T_2(s,t;A)={}&\frac{\zeta(s+t-1,A)}{s-1}
       -\frac12\zeta(s+t,A)\notag\\
 &+\sum_{k=1}^{M}\frac{B_{2k}}{(2k)!}(s)_{2k-1}
                      \zeta(s+t+2k-1,A),
\end{align}
and
\begin{align}\label{eq:EM3}
 T_3(s,t,v;A)={}&\frac{T_2(s+t-1,v;A)}{s-1}
       -\frac12T_2(s+t,v;A)\notag\\
 &+\sum_{k=1}^{M}\frac{B_{2k}}{(2k)!}(s)_{2k-1}
                      T_2(s+t+2k-1,v;A).
\end{align}
These follow by expanding the outer Hurwitz tail; the program deliberately
does not use the harmonic remainders. They are truncated asymptotic
approximations, not exact equalities to $Z_2,Z_3$.

For a ray $\bm c$, a circular sample gives
\[
 \frac1J\sum_{j=0}^{J-1}
 Z_3\bigl(1+c_1\rho e^{2\pi i j/J},
          1+c_2\rho e^{2\pi i j/J},
          1+c_3\rho e^{2\pi i j/J};a\bigr).
\]
The negative Laurent modes through order three cancel when $J>3$.
The remaining discrepancy includes positive-mode aliasing, quadrature
roundoff, and Euler--Maclaurin truncation. The recorded default tests use
$J=32$, $\rho=0.025$, $N=36$, $M=10$, and 55 working decimal digits,
with a conservative numerical comparison threshold of $10^{-25}$.
A separate calculation compares $N=36,M=10$ with $N=48,M=12$ at a
non-diagonal point in the ordinary convergence domain. No analytic
error enclosure for these numerical parameters is claimed. The proofs
in the body do not rely on their residuals.

\section{A finite implementation blueprint for arbitrary jets}
\label{app:algorithm}

For prescribed depth $d$, multi-index $\bm m$, and shift $a>0$, use the
following terminating recursion. Return $g_{m_1}(a)$ at depth one.
Otherwise compute the ordinary integral $A_{d;\bm m}$ from
\eqref{eq:Acoeff}; merge the first two indices into
$m_1+m_2+1$; multiply the recursive answer by
$\binom{m_1+m_2+1}{m_1+1}$; and add it. Symbolically this procedure
returns a finite linear combination of well-defined convergent
integrals and one Stieltjes coefficient. It is not a promise of an
automatic reduction of those integrals to classical constants.

For a straight-ray coefficient $[\eps^q]$, apply that procedure to the
finitely many multi-indices of degree $q+d-j$ appearing in each suffix
of \eqref{eq:allray}. For a nonlinear curve, first compute the prefix
valuations and the required truncations in Theorem \ref{thm:curve}, then
use reciprocal-series multiplication. This separates exact finite
algebra from special-function evaluation and from the analytic
normal-convergence proof.

\begin{thebibliography}{99}\small\raggedright
\setlength{\itemsep}{4pt}\setlength{\parskip}{0pt}
\bibitem{repo-target}
ProveIt Contributors, \emph{Audit, integration status, and further research},
\src{Analysis/Polylogarithms/docs/reports/stieltjes-harmonic-resonance/continuations/resonant-jets-and-shifts/sections/08-audit-research.tex}.
Git blob \src{b9c41fd8cec7205f60cc05c9f71e3873195659de}.
Repository: \url{https://github.com/VladimirReshetnikov/ProveIt}.

\bibitem{repo-double}
ProveIt Contributors, \emph{Directional and curved finite parts of double
Hurwitz zeta}, same continuation,
\src{sections/03-directional-hurwitz.tex}.
Git blob \src{cc9eebd46eeb6b62a4fa247eed41838133ec56fe}.

\bibitem{repo-readme}
ProveIt Contributors, \emph{Polylogarithms and their Arithmetic Bridges},
manuscript README, inspected 10 October 2026 (local-date research session).
\url{https://github.com/VladimirReshetnikov/ProveIt/tree/main/Analysis/Polylogarithms/docs/manuscript}.

\bibitem{repo-incoming}
ProveIt Contributors, incoming research directory, metadata inspected at
commit \src{0c740647c48baec7f0b4192723e977d69917dd52}.
\url{https://github.com/VladimirReshetnikov/ProveIt/tree/main/docs/incoming}.

\bibitem{mow}
K. Matsumoto, T. Onozuka and I. Wakabayashi,
\emph{Laurent series expansions of multiple zeta-functions of Euler--Zagier
type at integer points}, arXiv:1601.05918, especially Lemma 3.1 and
Theorems 1.3--1.4.
\url{https://arxiv.org/abs/1601.05918}.

\bibitem{gz}
L. Guo and B. Zhang, \emph{Renormalization of multiple zeta values},
Journal of Algebra \textbf{319} (2008), 3770--3809;
arXiv:math/0606076.
\url{https://arxiv.org/abs/math/0606076}.

\bibitem{gpxz}
L. Guo, S. Paycha, B. Xie and B. Zhang,
\emph{Double shuffle relations and renormalization of multiple zeta values},
arXiv:0906.0092.
\url{https://arxiv.org/abs/0906.0092}.

\bibitem{dlmf-zeta}
NIST Digital Library of Mathematical Functions, \S25.11,
\emph{Hurwitz Zeta Function}, especially the Mellin representation,
shift relation, and polygamma special values.
\url{https://dlmf.nist.gov/25.11}.

\bibitem{dlmf-polylog}
NIST Digital Library of Mathematical Functions, \S25.12,
\emph{Polylogarithms}.
\url{https://dlmf.nist.gov/25.12}.

\bibitem{dlmf-gamma}
NIST Digital Library of Mathematical Functions, \S\S5.5 and 5.7,
\emph{Functional Relations} and \emph{Series Expansions} for the Gamma
function.
\url{https://dlmf.nist.gov/5.5}; \url{https://dlmf.nist.gov/5.7}.
\end{thebibliography}
\end{document}
